% ============================================================================
% A Unified CT + MR Enterography Foundation-Model Pipeline for Crohn's Disease
% Lightweight modality conditioning matches dual specialists;
% domain-specific normalization causes negative transfer.
% ============================================================================
% Research Question:
% Can a single modality-conditioned foundation-model backbone extract
% radiological findings from BOTH CT and MR enterography as accurately as two
% independent modality specialists, and feed an MR-only survival head to
% predict the time to medical intervention in Crohn's disease?
% ============================================================================

\documentclass[preprint,12pt]{elsarticle}

\usepackage{amsmath,amssymb}
\usepackage{graphicx}
\usepackage{booktabs}
\usepackage{multirow}
\usepackage{xcolor}
\usepackage{hyperref}
\usepackage{algorithm}
\usepackage{algorithmic}
\usepackage{subcaption}
\usepackage{siunitx}

% Custom commands
\newcommand{\cindex}{C\text{-index}}
\newcommand{\todo}[1]{\textcolor{red}{\textbf{[TODO: #1]}}}

\journal{npj Digital Medicine}

\begin{document}

\begin{frontmatter}

\title{A Unified CT and MR Enterography Foundation-Model Pipeline for Crohn's Disease:\\
Lightweight Modality Conditioning Matches Dual Specialists, While Domain-Specific\\
Normalization Induces Negative Transfer}

\author[runi]{Ben Shaya}
\author[runi,szmc]{Moti Freiman\corref{cor1}}
\cortext[cor1]{Corresponding author}

\address[runi]{Faculty of Computer Science, Reichman University (RUNI), Herzliya, Israel}
\address[szmc]{Shaare Zedek Medical Center, Jerusalem, Israel}

\begin{abstract}
\textbf{Background:}
Crohn's disease (CD) is assessed with both MR enterography (MRE) and, frequently, CT enterography (CTE), yet deep-learning pipelines are almost always trained on a single modality. Whether one modality-conditioned foundation-model backbone can serve \textit{both} modalities as well as two dedicated specialists --- and whether such unification helps or hurts --- is unknown. We study this question in a two-stage prognostic pipeline that maps imaging to radiological findings and findings to time-to-event risk.

\textbf{Methods:}
We assembled 6{,}223 MRE and 1{,}830 CTE examinations from multiple Israeli centers, with patient-level (leakage-checked) splits and finding labels derived from radiology report text by a clinical information-extraction model. Stage~1 is a frozen DINOv3 backbone with LoRA adapters and dual-branch multi-instance learning (MIL) that detects 10 radiological findings. We compared \textbf{Option~B} (two independent CT and MR specialists) against \textbf{Option~C} (one unified backbone with modality conditioning), ablating three conditioning components: a modality token, feature-wise linear modulation (FiLM), and domain-specific normalization (DSBN). All models were warm-started from the same prior MR-only checkpoint, applied \emph{symmetrically} to both options. Stage~2 fits a penalized Cox model on the unified CT+MR cohort (6{,}011 studies, 4{,}191 MR + 1{,}820 CT, 5{,}256 patients) for surgery, steroid dependence, and biologic switch, with patient-grouped 5-fold cross-validation and per-modality C-index reporting. Evaluation used paired bootstrap confidence intervals, the integrated Brier score (IBS), leave-one-center-out (LOCO) validation, and decision-curve analysis (DCA).

\textbf{Results:}
The lead methodological finding is that the heaviest conditioning component, DSBN, causes \emph{statistically significant negative transfer}: removing it (token+FiLM vs.\ the full token+FiLM+DSBN model) improved MR macro-AUC by $+0.019$ (95\% CI $[+0.008, +0.030]$). With lightweight conditioning only (token+FiLM), the single unified model matched the MR specialist exactly (macro-AUC $0.849 = 0.849$) and was statistically equivalent to the CT specialist ($0.772$ vs.\ $0.760$; $+0.012$, 95\% CI $[-0.012, +0.035]$). Downstream, the same A2 backbone fed a \textbf{unified CT and MR Cox survival model} (6{,}011 studies, 4{,}191 MR + 1{,}820 CT, 817 surgery events): the combined model reached a surgery C-index of $0.816 \pm 0.015$ ($\Delta\cindex = +0.127$ over clinical features alone) with a CT-subset C-index of $0.767$, generalized across MR hospitals (LOCO surgery $0.805 \pm 0.023$), was well calibrated (IBS $= 0.072$), and achieved the highest net benefit in DCA across clinically relevant thresholds.

\textbf{Conclusion:}
A single modality-conditioned foundation-model backbone can replace two modality specialists for CD finding detection at no statistically significant cost --- halving the number of deployed models --- provided conditioning is kept lightweight; domain-specific normalization should be omitted because it over-parameterizes and degrades the dominant modality. We frame the unification as \emph{parity} (no detected accuracy difference, rather than a formally powered non-inferiority trial) rather than superiority; CT extends only the finding-detection stage, and all prognostic results are MR-based --- we make no claim that adding CT improves MR-based prognosis.
\end{abstract}

\begin{keyword}
Crohn's disease \sep CT enterography \sep MR enterography \sep Foundation models \sep Modality conditioning \sep Multi-instance learning \sep Survival analysis \sep Negative transfer
\end{keyword}

\end{frontmatter}

% ============================================================================
\section{Introduction}
\label{sec:introduction}
% ============================================================================

Crohn's disease (CD) is a chronic relapsing-remitting inflammatory bowel disease affecting over 3 million people in the United States, with rising global incidence \cite{torres2017crohns}. The clinical course is notoriously heterogeneous: some patients maintain long-term remission with minimal therapy, while others progress to surgical resection, steroid-dependent disease, or repeated escalation of biologic therapy \cite{peyrin2017lancet}. This unpredictability poses a fundamental challenge: \textit{which patients will deteriorate, and when?}

Cross-sectional enterography is central to CD assessment. MR enterography (MRE) offers radiation-free visualization of bowel-wall inflammation, thickening, strictures, fistulae, and extraintestinal complications \cite{defined_mri_cd}, and CT enterography (CTE) is widely used in acute and resource-constrained settings where MRE is unavailable. Validated activity scores such as MaRIA quantify disease activity at a single time point but do not predict future outcomes \cite{defined_maria}. The prognostic value of imaging --- its ability to predict \textit{what will happen to this patient} --- remains largely unexplored in an automated, quantitative framework, and almost all deep-learning work to date is single-modality.

\subsection{Gaps in Current Approaches}

\begin{enumerate}
    \item \textbf{One model per modality.} CT and MR pipelines are trained separately, doubling the engineering, validation, and deployment burden. Whether a single backbone can serve both modalities for CD --- and whether the field's default assumption that "shared learning across modalities helps" actually holds --- has not been tested for enterography. To our knowledge no published study trains a single deep model jointly on CTE and MRE for Crohn's.
    \item \textbf{Modality conditioning is under-characterized.} Recipes for serving multiple modalities from one backbone (modality tokens, FiLM, domain-specific normalization) are borrowed from segmentation and oncology, but their \textit{relative} contribution --- and the risk that heavier conditioning hurts the dominant modality (\emph{negative transfer}: added cross-modal capacity that lowers rather than raises performance) --- has not been ablated in this setting.
    \item \textbf{Classification is not prediction.} Most deep learning in CD addresses binary activity or scoring at a single time point \cite{dl_ibd_review}, answering "is the patient sick now?" rather than "will this patient need surgery, and when?" --- the more actionable question.
    \item \textbf{Venue-grade evidence is often missing.} External (multi-center) validation, calibration reporting (Brier score), and clinical net-benefit analysis are expected for digital-medicine deployment claims but are frequently absent in CD imaging studies.
\end{enumerate}

\subsection{Research Question}

\begin{quote}
\textit{Can a single modality-conditioned foundation-model backbone extract radiological findings from both CT and MR enterography as accurately as two independent modality specialists, and feed a survival head that predicts the time to medical intervention in Crohn's disease?}
\end{quote}

We address this with a two-stage pipeline: a frozen DINOv3 backbone with LoRA adapters and dual-branch multi-instance learning (MIL) detects 10 radiological findings (Stage~1); a penalized Cox model maps findings plus clinical features to individualized time-to-event predictions (Stage~2). We directly compare two independent specialists (Option~B) against one unified, modality-conditioned backbone (Option~C), and we ablate which conditioning components help or harm.

\subsection{Contributions}

\begin{enumerate}
    \item \textbf{Domain-specific normalization causes significant negative transfer (lead finding).} In a controlled ablation ladder over modality conditioning (token, FiLM, DSBN), removing DSBN improves MR macro-AUC by $+0.019$ (95\% CI $[+0.008, +0.030]$, paired bootstrap). Lightweight conditioning (modality token + FiLM) is the sweet spot; heavier conditioning over-parameterizes and degrades the dominant modality.

    \item \textbf{One unified model matches two specialists (non-inferiority).} With token+FiLM conditioning, a single backbone matches the MR specialist exactly (macro-AUC $0.849$) and is statistically equivalent to the CT specialist ($0.772$ vs.\ $0.760$), and matches the specialists' downstream survival performance. Because the same warm-start is applied \emph{symmetrically} to both options, the head-to-head is fair by construction. This halves the number of models to validate and maintain at no statistically significant cost.

    \item \textbf{An MR-based prognostic pipeline with venue-grade evidence.} The unified model's survival head predicts surgery (C-index $0.825$; $\Delta\cindex = +0.145$ over clinical-only), steroid dependence, and biologic switch, with external validation (LOCO surgery $0.799 \pm 0.027$), calibration (IBS $0.085$), and decision-curve net benefit across clinically relevant thresholds.

    \item \textbf{A reproducible CT-localization pipeline.} We adapt an organ-segmentation crop (TotalSegmentator $\rightarrow$ small-bowel + colon $\rightarrow$ bounding-box crop $\rightarrow$ 16 slices) so that CTE is bowel-localized and consumable by the same MIL backbone as MRE.

    \item \textbf{A unified CT+MR survival model.} Updated clinical outcomes link time-to-event labels to 1{,}820 CT studies, enabling --- to our knowledge for the first time --- a Cox prognostic model jointly fit on CT and MR enterography. CT contributes a clinically meaningful surgery C-index of $0.767$ on its own subset, and slightly exceeds MR for biologic switch ($0.691$ vs.\ $0.655$).
\end{enumerate}


% ============================================================================
\section{Related Work}
\label{sec:related}
% ============================================================================

\subsection{Deep Learning for Crohn's Disease Imaging}

Deep learning has been applied to gastrointestinal imaging, including endoscopy-based IBD assessment and MRI-based quantification of CD activity. Rimola et al.\ established the MaRIA score as a validated quantitative measure of activity from MRI \cite{defined_maria}, and subsequent work automated scoring with convolutional networks \cite{defined_mri_cd}; reviews highlight AI for bowel-wall segmentation, activity classification, and treatment-response prediction \cite{dl_ibd_review}. CTE deep learning has separately reached strong activity/fibrosis performance in single-modality cohorts. Critically, the field splits cleanly by modality: to our knowledge no prior study trains a single deep model jointly on CTE and MRE for Crohn's, making a unified enterography model a publishable gap.

\subsection{Foundation Models Across CT and MR}

The dominant 2024--2026 paradigm is one frozen foundation backbone with modality-aware adapters. Frozen DINOv2/DINOv3 features generalize across X-ray, CT, and MR benchmarks with a single encoder \cite{baharoon2024dinov2}, and purpose-built cross-modal backbones trained on mixed CT and MR (e.g., Curia) degrade far less on out-of-distribution MR than domain-general embeddings \cite{curia2025}. Embedding-only fusion of frozen foundation features has reached state-of-the-art oncologic prognosis at scale \cite{honeybee2025}. We adopt this paradigm: a frozen DINOv3 backbone \cite{dinov3} with LoRA adapters \cite{hu2021lora}, shared across CT and MR.

\subsection{Modality Conditioning}

Several lightweight techniques let one backbone serve multiple modalities: a learned \textit{modality token}, \textit{feature-wise linear modulation} (FiLM), which predicts per-channel scale/shift parameters \cite{perez2018film}, and \textit{domain-specific normalization} (DSBN), which keeps separate normalization statistics per domain \cite{chang2019dsbn}. These have shown gains in cross-modal segmentation and adaptation, but their relative contribution --- and the possibility that the heaviest component induces negative transfer on the dominant modality --- has not been characterized for enterography. We ablate all three.

\subsection{Multi-Instance Learning and Survival from Images}

Attention-based MIL \cite{ilse2018attention} weights instances (slices) under bag-level (study-level) labels and has succeeded in pathology \cite{campanella2019clinical}, chest CT \cite{li2020dual}, and brain MRI \cite{defined_mil_brain}. Time-to-event prediction from images spans end-to-end deep survival (DeepSurv \cite{katzman2018deepsurv}, DeepHit \cite{lee2018deephit}) and two-stage approaches that feed learned features to a classical survival model. Recent npj Digital Medicine multimodal-survival pipelines (e.g., multimodal transformers with Cox/AFT heads and external validation \cite{gliosurv2025}, deep multimodal fusion with feature-category ablations \cite{cervipro2025}, and patho-radiomic fusion \cite{prismcrc2025}) set the venue norm of $\geq 2$-cohort external validation with calibration reporting. We adopt the two-stage design for supervision and statistical-efficiency reasons (Section~\ref{sec:stage2}) and report the venue-grade evidence those works establish.


% ============================================================================
\section{Materials and Methods}
\label{sec:methods}
% ============================================================================

Figure~\ref{fig:architecture} overviews the two-stage pipeline and the Option~B vs.\ Option~C comparison.

\subsection{Study Population and Data}
\label{sec:data}

This retrospective study pooled CD examinations from multiple Israeli centers (Mor, Shaare Zedek Medical Center [SZMC], Assuta; with additional report-derived labels from Rambam). The study was approved by the institutional review board. Original imaging and clinical databases were treated as read-only.

\subsubsection{Imaging Data}
Stage~1 used \textbf{6{,}223 MRE} examinations (reused from the prior MR-only cohort to enable a valid warm-start) and \textbf{1{,}830 CTE} examinations ($\approx 23\%$ of the imaging cohort), the latter from SZMC and Assuta. Each study was reduced to 16 axial slices resized to $224 \times 224$. For MRE, T2-weighted and T1-weighted post-contrast sequences were used; CTE was routed through the T2 branch with the T1 branch zeroed (Section~\ref{sec:stage1}).

\paragraph{CT bowel localization.}
Because CT spans the whole abdomen and thorax, we localized the bowel before slice sampling: TotalSegmentator \cite{wasserthal2023totalseg} (fast mode) produced small-bowel and colon masks, whose union bounding box (with a 10-voxel pad) cropped the volume; 16 slices were then evenly sampled on the cropped volume. This mirrors the MRE bowel-crop pipeline, so localization comes from the crop rather than the slice sampler.

\subsubsection{Radiological Finding Labels}
Ten radiological findings (Table~\ref{tab:labels}) were extracted from the free-text radiology reports (Hebrew enterography reports) by a clinical information-extraction model, applied identically to MR and CT reports. We emphasize that these labels are \emph{model-derived from report text, not prospective human ground truth}; both Option~B and Option~C are evaluated against the \emph{same} labels, so the comparison between them is unaffected by label noise. Graded inflammation outputs were binarized at grade~$\geq 1$; on the overlapping cohort this reproduced the prior validated binary labels exactly (6{,}223/6{,}223 studies), supporting the mapping. \todo{Report the HSMP-BERT extraction precision/recall against a radiologist-annotated subset --- all \emph{absolute} metrics inherit this label ceiling, although the B-vs-C comparison does not.}

\begin{table}[h]
\centering
\caption{Radiological finding labels with prevalence in the MR training set.}
\label{tab:labels}
\begin{tabular}{@{}llcc@{}}
\toprule
& \textbf{Finding} & \textbf{Prevalence (\%)} & \textbf{Tier} \\
\midrule
1 & Ileal inflammation        & 33.1\% & 1 \\
2 & Ileal wall enhancement    & 40.4\% & 1 \\
3 & Ileal wall thickening     & 42.9\% & 1 \\
4 & Restricted diffusion (DWI)& 26.8\% & 1 \\
5 & Ileal stenosis            & 22.0\% & 1 \\
6 & Pre-stenotic dilatation   & 14.7\% & 1 \\
7 & Comb sign                 & 12.1\% & 1 \\
8 & Fistula                   & 6.6\% & 2 \\
9 & Colonic inflammation      & 6.3\% & 2 \\
10 & Mesenteric edema          & 4.6\% & 2 \\
\bottomrule
\end{tabular}
\end{table}

\subsubsection{Survival Outcomes}
\label{sec:survival_data}
An updated clinical-outcomes release linked time-to-event labels for 1{,}820 CT studies in our cohort, allowing \textbf{Stage~2 to be unified across CT and MR}. Stage~2 used \textbf{6{,}011 examinations} (4{,}191 MR + 1{,}820 CT) from 5{,}256 unique patients, with time zero set to the imaging date (to avoid immortal-time bias). Three outcomes were defined:
\begin{itemize}
    \item \textbf{Surgical intervention} --- 817 events (687 MR + 130 CT).
    \item \textbf{Steroid dependence} --- 292 events (226 MR + 66 CT).
    \item \textbf{Biologic therapy switch} --- 1{,}232 events (968 MR + 264 CT).
\end{itemize}
The event indicator was derived from the post-imaging signed time-to-event ($>0 \Rightarrow$ event), non-events were right-censored at a fixed maximum follow-up of 10 years, and the same convention was applied to MR and CT. Events-per-variable on 23 features is 35 (surgery), 13 (steroid), and 54 (biologic), well above the conventional minimum of 10 for stable Cox fitting.

\subsubsection{Clinical Features}
Thirteen clinical features (6 continuous, 7 binary) were extracted from the clinical outcomes table, representing disease characteristics and treatment history at the time of MRI: time from diagnosis to MRI, age at diagnosis, Charlson Comorbidity Index, disease-activity cluster, socioeconomic status, diagnostic delay; and prior surgery, prior steroid dependence, prior biologic switch, prior extraintestinal manifestations, prior perianal disease, sex, and a clinical-data-availability indicator. Missing values were imputed from the training set only (median/mode), with a \texttt{has\_clinical} indicator to allow the model to discount imputed values.

\subsubsection{Data Splitting}
All splits were \textit{patient-level} (a patient appears in only one split; leakage explicitly checked across 7{,}395 patients). Stage~1 was evaluated on held-out test sets of $n = 628$ MR and $n = 209$ CT studies. We note that the patient-level resplit of the new cohort differs slightly from the prior split, so test sets are not identical to earlier phases.


% ============================================================================
\subsection{Stage 1: Unified Imaging Biomarker Extraction}
\label{sec:stage1}
% ============================================================================

Stage~1 transforms raw slices into 10 sigmoid-calibrated finding probabilities that serve as Stage~2 inputs.

\subsubsection{Backbone and MIL Head}
A \textbf{frozen DINOv3-Base} backbone \cite{dinov3} extracts per-slice features; \textbf{LoRA adapters} \cite{hu2021lora} ($\approx 8.5\%$ of parameters trainable) adapt it cheaply and stably. A dual-branch design processes T2 and T1 (for MR; CT uses the T2 branch with T1 zeroed via the existing missing-modality mask), projecting to a fused per-slice representation. Ten \emph{independent} gated-attention heads \cite{ilse2018attention}, one per finding, pool slices into bag representations:
\begin{equation}
    a_k^{(i)} = \text{softmax}\!\left( \mathbf{w}^{(i)\top} \!\left( \tanh(\mathbf{V}^{(i)} \mathbf{h}_k) \odot \sigma(\mathbf{U}^{(i)} \mathbf{h}_k) \right) \right), \qquad
    \mathbf{z}^{(i)} = \sum_{k=1}^{16} a_k^{(i)} \mathbf{h}_k ,
\end{equation}
where $\mathbf{h}_k$ is the fused representation of slice $k$ and $i$ indexes the finding. Each bag representation is concatenated with clinical features and passed through a per-finding MLP.

\subsubsection{Option B vs.\ Option C, and Modality Conditioning}
We compare two designs:
\begin{itemize}
    \item \textbf{Option~B (specialists):} two independent backbones, B-MR (trained on MR) and B-CT (trained on CT), with no conditioning.
    \item \textbf{Option~C (unified):} one backbone trained jointly on CT+MR (WeightedRandomSampler balancing modality$\times$label), with modality conditioning. We ablate three components:
    \begin{itemize}
        \item \textbf{Modality token} --- a learned embedding indicating CT vs.\ MR;
        \item \textbf{FiLM} \cite{perez2018film} --- per-channel scale/shift on the backbone features, predicted from the modality token;
        \item \textbf{DSBN} \cite{chang2019dsbn} --- modality-specific normalization on the fused representation.
    \end{itemize}
\end{itemize}
The ablation ladder is: \textbf{A1} (token only), \textbf{A2} (token+FiLM), \textbf{A3} (token+FiLM+DSBN, the "full" model), and \textbf{A5} (DSBN, MR-only control). The labels index a larger experiment matrix (hence the non-contiguous numbering); we report the informative rungs.

\subsubsection{Symmetric Warm-Start (fair comparison)}
All Stage~1 models --- both specialists and unified variants --- were warm-started from the \emph{same} prior MR-only DINOv3 checkpoint. Applying the warm-start \emph{symmetrically} to Option~B and Option~C is what makes the head-to-head fair: any shared initialization advantage is conferred equally on both sides, so the comparison isolates the effect of modality conditioning, not initialization.

\subsubsection{Training Configuration}
Models were trained on a single NVIDIA A6000 (48~GB) in bfloat16 with a tier-aware asymmetric loss \cite{ridnik2021asl} for class imbalance, AdamW with differential learning rates, EMA, and early stopping. Warm-starting from the prior optimum, the unified and specialist models converged within a few epochs.


% ============================================================================
\subsection{Stage 2: Time-to-Event Prediction}
\label{sec:stage2}
% ============================================================================

For each MR study, the trained Stage~1 model produces 10 probabilities $\hat{p}_1,\ldots,\hat{p}_{10}$, concatenated with 13 clinical features to form a 23-dimensional vector. A separate elastic-net penalized Cox model \cite{polsterl2020scikit} is fit per outcome:
\begin{equation}
    h(t \mid \mathbf{x}) = h_0(t)\exp(\boldsymbol{\beta}^\top \mathbf{x}),
\end{equation}
with the Breslow estimator giving the baseline cumulative hazard and individualized survival curves $\hat{S}(t\mid\mathbf{x})$. The two-stage design is motivated by (a) the supervision advantage (Stage~1 sees far more labeled studies than there are survival events), (b) interpretability of the intermediate biomarkers, and (c) statistical efficiency: at moderate event counts, a low-dimensional Cox model respects the events-per-variable (EPV) constraint (surgery: 659 events / 23 features $\approx 29$, well above the conventional minimum of 10) where an end-to-end deep survival model would overfit.

\subsubsection{Evaluation Protocol}
We report: the concordance index (C-index; the probability that the model assigns higher risk to the patient with the sooner event, $0.5=$ chance) with bootstrap ($n=1{,}000$) confidence intervals; an ablation over feature sets (combined 23-d vs.\ imaging-only 10-d vs.\ clinical-only 13-d), with $\Delta\cindex$ quantifying the added value of imaging; the integrated Brier score (IBS; squared survival-prediction error integrated over time, lower is better) and time-dependent AUC for calibration; \textbf{leave-one-center-out} (LOCO) external validation across the three MR centers; and decision-curve analysis (DCA) for clinical net benefit. For Stage~1, model pairs (B vs.\ C, and ablation rungs) were compared with paired bootstrap confidence intervals on the test-set macro-AUC (the unweighted mean of the 10 per-finding AUCs).


% ============================================================================
\section{Results}
\label{sec:results}
% ============================================================================

\subsection{Stage 1: Unified vs.\ Specialist, and the Cost of DSBN}

Table~\ref{tab:ladder} reports test-set macro-AUC for the specialists and the unified ablation ladder.

\begin{table}[h]
\centering
\caption{Stage~1 finding-detection macro-AUC on the held-out test set (MR $n=628$, CT $n=209$). All models share a symmetric warm-start. Higher is better.}
\label{tab:ladder}
\begin{tabular}{@{}llcc@{}}
\toprule
\textbf{Config} & \textbf{Conditioning} & \textbf{MR} & \textbf{CT} \\
\midrule
B-MR / B-CT (specialists) & ---                       & \textbf{0.849} & 0.760 \\
A1 (unified)              & token                     & 0.843 & \textbf{0.776} \\
\textbf{A2 (unified)}$^\star$ & token + FiLM          & \textbf{0.849} & 0.772 \\
A5 (control)             & DSBN, MR-only             & 0.835 & --- \\
A3 (unified, ``full'')    & token + FiLM + DSBN       & 0.830 & 0.756 \\
\bottomrule
\end{tabular}
\end{table}

\paragraph{Per-finding A2 performance.}
Table~\ref{tab:per_finding_a2} reports per-finding AUC and Youden-optimal sensitivity/specificity for A2 on the MR and CT test sets. MR macro-AUC $= 0.849$ ($n=628$); CT macro-AUC $= 0.772$ ($n=209$, with DWI not assessed on CT).

\begin{table}[h]
\centering
\caption{Per-finding test-set AUC and Youden-optimal sensitivity/specificity for the unified A2 model. ``---'' denotes findings not assessable on CT (no DWI sequence in CTE).}
\label{tab:per_finding_a2}
\small
\begin{tabular}{@{}lcccccc@{}}
\toprule
 & \multicolumn{3}{c}{\textbf{MR} ($n=628$)} & \multicolumn{3}{c}{\textbf{CT} ($n=209$)} \\
\cmidrule(lr){2-4}\cmidrule(lr){5-7}
\textbf{Finding} & \textbf{AUC} & \textbf{Sens.} & \textbf{Spec.} & \textbf{AUC} & \textbf{Sens.} & \textbf{Spec.} \\
\midrule
Ileal inflammation        & 0.821 & 0.745 & 0.750 & 0.755 & 0.574 & 0.843 \\
Ileal wall enhancement    & 0.866 & 0.775 & 0.798 & 0.735 & 0.744 & 0.634 \\
Ileal wall thickening     & 0.865 & 0.773 & 0.799 & 0.758 & 0.542 & 0.912 \\
Restricted diffusion (DWI)& 0.841 & 0.818 & 0.739 & ---   & ---   & ---   \\
Ileal stenosis            & 0.844 & 0.733 & 0.813 & 0.834 & 0.898 & 0.681 \\
Pre-stenotic dilatation   & 0.852 & 0.752 & 0.803 & 0.845 & 0.868 & 0.778 \\
Comb sign                 & 0.806 & 0.829 & 0.658 & 0.784 & 0.667 & 0.872 \\
Fistula                   & 0.852 & 0.927 & 0.684 & 0.794 & 0.857 & 0.678 \\
Colonic inflammation      & 0.843 & 0.806 & 0.811 & 0.768 & 0.536 & 0.890 \\
Mesenteric edema          & 0.896 & 0.810 & 0.842 & 0.677 & 0.667 & 0.711 \\
\midrule
\textbf{Macro-AUC}        & \textbf{0.849} & --- & --- & \textbf{0.772} & --- & --- \\
\bottomrule
\end{tabular}
\end{table}

\paragraph{DSBN causes significant negative transfer.}
Removing DSBN (A2 vs.\ A3) improved MR macro-AUC by $+0.019$ (95\% CI $[+0.008, +0.030]$, paired bootstrap) --- the confidence interval excludes zero. DSBN is the only conditioning component that significantly hurts, and it hurts the dominant modality. The MR-only DSBN control (A5, $0.835$) likewise underperforms token+FiLM, confirming the effect is attributable to DSBN rather than to joint training.

\paragraph{Lightweight conditioning matches the specialists (non-inferiority).}
The token+FiLM unified model (A2) \emph{matches the MR specialist exactly} ($0.849 = 0.849$) and is statistically equivalent to the CT specialist on CT ($0.772$ vs.\ $0.760$; $+0.012$, 95\% CI $[-0.012, +0.035]$; the interval includes zero, so the CT difference is a non-significant trend on a small test set, $n=209$). Token-only conditioning (A1) is also statistically indistinguishable from both specialists. Thus a single unified backbone with lightweight conditioning replaces two specialists at no statistically significant cost.

\paragraph{Nested cross-validation: the ranking is stable across patient-level CV folds.}
To verify that the production choice of A2 is not an artifact of a single train/val/test split, we re-trained the entire conditioning ladder (A1, A2, A3) on five patient-level CV folds and evaluated each (fold, variant) pair on the fold's held-out test set. Table~\ref{tab:nested_cv} reports per-variant mean $\pm$ SD across the five folds. On the MR test sets, A2 outperforms A3 in \emph{every one of the five folds} (paired-fold mean $\Delta = +0.0075$, 95\% CI $[+0.003, +0.012]$, CI excludes zero), confirming the DSBN negative-transfer finding as a stable property of the conditioning ladder rather than a single-split artifact. On the CT test sets, A2 and A3 are statistically indistinguishable (paired-fold mean $\Delta \approx 0$, large fold-to-fold variance), so the small CT trend in favor of A3 visible in Table~\ref{tab:ladder} is within noise.

\begin{table}[h]
\centering
\caption{Nested-CV test-set macro-AUC across 5 patient-level folds (mean $\pm$ SD). The conditioning ladder was retrained from scratch on each fold's train+val split. A2 achieves the best MR mean with the lowest fold-to-fold SD; A3 (with DSBN) consistently underperforms on MR.}
\label{tab:nested_cv}
\begin{tabular}{@{}lccc@{}}
\toprule
 & \textbf{A1 (token)} & \textbf{A2 (token+FiLM)}$^\star$ & \textbf{A3 (token+FiLM+DSBN)} \\
\midrule
\textbf{MR} & $0.896 \pm 0.010$ & $\mathbf{0.897 \pm 0.007}$ & $0.889 \pm 0.007$ \\
\textbf{CT} & $0.758 \pm 0.024$ & $0.768 \pm 0.024$ & $0.769 \pm 0.016$ \\
\bottomrule
\end{tabular}
\end{table}

\paragraph{Leave-one-modality-out: CT supervision neither helps nor hurts MR.}
To probe whether the unified model's MR performance is driven by joint CT+MR training or simply by the shared backbone and warm-start, we trained an additional token+FiLM model on MR-only and evaluated both on the held-out MR test set ($n=628$). The joint A2 (CT+MR) and the MR-only variant achieved essentially identical macro-AUC ($0.8486$ vs.\ $0.8484$; $\Delta = +0.0002$, paired-bootstrap 95\% CI $[-0.0086, +0.0090]$; the joint model was higher in $51\%$ of bootstrap resamples). CT supervision therefore has \emph{no detectable effect} on MR finding detection in either direction --- supporting our explicit non-claim that adding CT improves MR-based prediction, while showing that joint training does \emph{not} degrade the dominant modality. The unified model's value is operational (one model serving both modalities), not informational.

\subsection{Stage 2: Unified CT and MR Time-to-Event Prediction}

Using the production unified model (A2) to generate imaging biomarkers, we fit a penalized Cox model on the unified 6{,}011-study cohort (4{,}191 MR + 1{,}820 CT) with patient-grouped 5-fold cross-validation. For each held-out fold we report the C-index on the combined test cohort and, separately, on the MR and CT subsets of that fold. To our knowledge this is the first time a Crohn's-disease prognostic Cox model has been fit jointly on CT and MR enterography.

\begin{table}[h]
\centering
\caption{Stage~2 C-index (5-fold patient-grouped CV, mean $\pm$ SD) for the unified A2 model, on the combined test cohort and on its MR and CT subsets. $\Delta$MRI = combined $-$ clinical-only on the matching subset.}
\label{tab:ablation}
\small
\begin{tabular}{@{}lccc|ccc|c@{}}
\toprule
 & \multicolumn{3}{c|}{\textbf{Combined cohort}} & \multicolumn{3}{c|}{\textbf{Per-modality subset}} & \\
\textbf{Outcome} & \textbf{Combined} & \textbf{Imaging} & \textbf{Clinical} & \textbf{MR} & \textbf{CT} & \textbf{Clin. (MR/CT)} & \textbf{$\Delta$MRI} \\
\midrule
Surgery   & $\mathbf{0.816 \pm 0.015}$ & $0.776 \pm 0.019$ & $0.689 \pm 0.018$ & $0.815 \pm 0.012$ & $0.767 \pm 0.050$ & $0.670 / 0.640$ & $+0.127$ \\
Steroid   & $0.644 \pm 0.021$ & $0.534 \pm 0.028$ & $0.646 \pm 0.033$ & $0.663 \pm 0.028$ & $0.559 \pm 0.062$ & $0.653 / 0.581$ & $-0.002$ \\
Biologic  & $0.671 \pm 0.018$ & $0.600 \pm 0.018$ & $0.627 \pm 0.017$ & $0.655 \pm 0.016$ & $0.691 \pm 0.047$ & $0.616 / 0.600$ & $+0.044$ \\
\bottomrule
\end{tabular}
\end{table}

\paragraph{Imaging contributes prognostically on both modalities.}
For surgery, imaging adds $\Delta\cindex = +0.127$ over clinical-only on the combined cohort, and the per-modality split shows CT delivers a meaningful prognostic signal in its own right (CT-only C-index $0.767$, vs.\ $0.640$ from clinical features alone on the CT subset; $\Delta = +0.127$). For biologic switch, CT in fact \emph{exceeds} MR in the per-modality C-index ($0.691$ vs.\ $0.655$), a finding that would have been invisible under the previous MR-only Stage~2.

\paragraph{Adding CT does not degrade MR.}
The MR-subset surgery C-index under the unified Cox ($0.815 \pm 0.012$) is statistically indistinguishable from the corresponding MR-only Cox fit on the prior cohort ($0.825 \pm 0.020$; difference within fold-to-fold SD). Joint training of Cox on CT+MR is therefore safe for the MR side, consistent with the Stage-1 leave-one-modality-out finding.

\paragraph{Steroid dependence remains clinically driven.}
Steroid dependence is dominated by clinical history on both modalities; imaging-only is near chance, mirroring our previous result.

\subsection{External Validation, Calibration, and Clinical Utility}

\paragraph{Multi-center generalization (LOCO).}
Leaving out each MR center in turn (Mor / SZMC / Assuta) and refitting Cox per fold on the unified cohort, the held-out-center surgery C-index was $\mathbf{0.805 \pm 0.023}$ (combined model), $0.744 \pm 0.033$ (imaging-only), and $0.659 \pm 0.045$ (clinical-only). Imaging retains its added value across held-out centers ($\Delta\cindex = +0.146$ vs.\ clinical-only under LOCO). CT-only LOCO is not feasible because center labels are unavailable for the CT cohort in this release; we report MR-LOCO as the external-generalization proxy.

\paragraph{Calibration.}
The combined surgery model achieved an integrated Brier score of $\mathbf{0.072}$ on the unified CT+MR cohort --- the lowest (best) of all configurations and better than imaging-only (0.077) or clinical-only (0.085) --- with time-dependent AUC of $\approx 0.85$ across the 1/2/3-year horizons for surgery (Figure~\ref{fig:calibration}).

\paragraph{Decision-curve analysis.}
Across the clinically relevant threshold range ($\approx 0.02$--$0.60$), the combined model on the unified cohort achieved the highest net benefit relative to treat-all and treat-none defaults (Figure~\ref{fig:dca}), with a mean net-benefit gain of $+0.020$ over the clinical-only Cox at the 5-year surgery horizon ($n=6{,}011$, $817$ events), supporting clinical actionability.

\subsection{Feature Importance and Risk Stratification}

Consistent with the surgery ablation, the imaging biomarkers corresponding to penetrating and stricturing phenotypes (fistula, stenosis, wall thickening, wall enhancement) carry the largest standardized surgical hazards on the unified CT+MR cohort, with seven of eleven elastic-net-selected features being MRI-derived (Table~\ref{tab:hazard_ratios}). The clinical \texttt{prior surgery} indicator is strongly protective ($\text{HR} = 0.62$), reflecting that patients who already underwent surgery are at lower risk of a \emph{subsequent} surgical event in the follow-up window. Predicted-risk tertiles separate event curves with high significance on the unified cohort (Figure~\ref{fig:km_curves}).

\begin{table}[h]
\centering
\caption{Non-zero elastic-net Cox coefficients and hazard ratios for surgery on the unified CT+MR A2 cohort ($n = 6{,}011$, $817$ events; $11$ of $23$ candidate features retained at the chosen regularization $\alpha = 0.033$). Features were standardized; HRs are reported per one standard deviation. Seven of eleven selected features are MRI-derived.}
\label{tab:hazard_ratios}
\small
\begin{tabular}{@{}lccc@{}}
\toprule
\textbf{Feature} & \textbf{Source} & \textbf{HR per 1 SD} & \textbf{Direction} \\
\midrule
Fistula                       & MRI (MIL) & 1.363 & Risk $\uparrow$ \\
Ileal stenosis                & MRI (MIL) & 1.290 & Risk $\uparrow$ \\
Ileal wall thickening         & MRI (MIL) & 1.224 & Risk $\uparrow$ \\
Ileal wall enhancement        & MRI (MIL) & 1.129 & Risk $\uparrow$ \\
Colonic inflammation          & MRI (MIL) & 1.108 & Risk $\uparrow$ \\
Restricted diffusion (DWI)    & MRI (MIL) & 1.096 & Risk $\uparrow$ \\
Pre-stenotic dilatation       & MRI (MIL) & 1.012 & Risk $\uparrow$ \\
Prior steroid dependence      & Clinical  & 1.018 & Risk $\uparrow$ \\
Socioeconomic status (SES)    & Clinical  & 0.970 & Risk $\downarrow$ \\
Sex (male)                    & Clinical  & 0.973 & Risk $\downarrow$ \\
Prior surgery                 & Clinical  & 0.622 & Risk $\downarrow$ \\
\bottomrule
\end{tabular}
\end{table}


% ============================================================================
\section{Discussion}
\label{sec:discussion}
% ============================================================================

We asked whether one modality-conditioned foundation-model backbone can serve both CT and MR enterography as well as two specialists, within a two-stage prognostic pipeline. The answer is a qualified yes --- \emph{provided conditioning is kept lightweight}.

\subsection{Domain-Specific Normalization Over-Parameterizes and Hurts}

Our cleanest, most generalizable result is methodological: among modality-conditioning components, DSBN is the one that significantly degrades performance, and it degrades the dominant modality (MR $+0.019$ macro-AUC recovered by removing it, CI excluding zero). A modality token and FiLM are nearly free and harmless-to-helpful; DSBN's separate per-modality normalization adds parameters that, at this scale, the joint task cannot support --- classic negative transfer. The practical recommendation is concrete: \textbf{condition with a token and FiLM, and omit DSBN}. This also reconciles an apparent contradiction with the cross-modal literature, where DSBN helps in large-scale adaptation: at a moderate, imbalanced CT+MR scale ($\approx 23\%$ CT), the heavier component flips from helpful to harmful.

\subsection{Unification as Parity, Not Superiority}

We deliberately frame Option~C as achieving \emph{parity} with Option~B rather than superiority. With token+FiLM, the unified model matches the MR specialist exactly and is statistically equivalent to the CT specialist (CT trend favorable but underpowered, $n=209$, CI includes zero), and matches/edges the specialists downstream. We use ``parity'' in the descriptive sense of no detected difference: this is \emph{not} a pre-specified, powered non-inferiority test with a clinical margin, and the CT comparison in particular is an underpowered null. The deployment claim we make --- one model rather than two, at no detected accuracy cost --- is robust to this caveat; a formal non-inferiority margin would strengthen it. The deployment value is real and modest: one model instead of two to train, validate, calibrate, and maintain, with no statistically significant accuracy cost. The comparison is fair because the warm-start is applied symmetrically to both options; a cold-start experiment would instead probe the \emph{mechanism} of transfer (which we do not claim) rather than the deployment question (which we do). We therefore avoid any \emph{Stage-1} "CT improves MR" claim: the leave-one-modality-out experiment shows CT supervision has no detectable effect on MR finding detection. The Stage-2 unified Cox is a separate, complementary claim --- CT contributes a prognostic signal of its own ($n=1{,}820$, surgery $C = 0.767$) and adding CT to the unified Cox does not degrade MR ($C = 0.815$ on the MR subset, indistinguishable from the prior MR-only Cox).

\subsection{MRI Reveals Prognostic Information Beyond Clinical History}

For surgery, imaging contributes substantial prognostic value ($\Delta\cindex = +0.145$ over clinical-only), with penetrating/stricturing imaging phenotypes carrying the strongest hazard. A patient with radiological evidence of fistula, stenosis, or pre-stenotic dilatation faces elevated surgical risk even when clinical history appears unremarkable. Outcome-specific pathways are evident: steroid dependence is driven by treatment history (a clinical phenomenon), while biologic switch reflects both morphology and history.

\subsection{Stage 2 Plateaus Across Backbones --- an EPV Hypothesis}

Stage~2 performance is remarkably stable across backbones and conditioning choices (the unified, full-DSBN, and specialist variants all land near a surgery C-index of $\approx 0.82$). We interpret this cautiously. With the unified cohort, EPV at 23 features is comfortable ($\approx 35$ for surgery, $13$ for steroid, $54$ for biologic), so this is not a raw degrees-of-freedom problem; rather, the 10-finding $\rightarrow$ Cox \emph{interface} may already saturate the prognostic signal extractable at this event count, so small Stage-1 macro-AUC gains need not perturb the finding probabilities enough to move downstream risk. We present this as a hypothesis, not a tested claim; confirming it would require manipulating event count or the survival-head dimensionality directly.

\subsection{Clinical Utility}

External (LOCO) generalization, low IBS, and favorable DCA net benefit collectively meet the digital-medicine evidence bar for a deployment-oriented prognostic tool, while remaining honest about scope.

\subsection{Limitations}

\begin{itemize}
    \item \textbf{Labels are model-derived.} The 10 findings come from a report-text information-extraction model, not prospective human annotation. Both options are scored against the same labels, so the \emph{comparison} is unaffected, but absolute performance inherits label noise.
    \item \textbf{Single internal cohort; multi-center external validation pending.} Although A2's edge over A3 is confirmed by nested cross-validation across five patient-level CV folds (Table~\ref{tab:nested_cv}), and the leave-one-modality-out experiment shows the joint-vs.\ MR-only choice is not load-bearing, all evaluation is internal to the contributing Israeli centers. True external validation on an independent cohort would strengthen the deployment claim.
    \item \textbf{CT Stage-1 comparison is underpowered.} The CT Stage-1 test set is small ($n=209$), so the CT specialist-vs-unified comparison is a trend rather than a powered result. Stage 2 now includes CT (1{,}820 studies, 130 surgery events) thanks to an updated outcomes release, but center labels are not yet available for CT, so multi-center LOCO is reported on MR only.
    \item \textbf{Retrospective design and proportional-hazards assumption.} Selection biases cannot be fully excluded; time-varying effects were not modeled.
    \item \textbf{Per-study modeling.} Stage~2 operates at the study level under patient-level splits; per-patient temporal aggregation of multi-scan patients is left to future work.
\end{itemize}

\subsection{Future Directions}

\begin{enumerate}
    \item \textbf{Leave-one-modality-out.} Train token+FiLM on MR-only and compare to the joint CT+MR model on the MR test set; only if the joint model is non-inferior or better on MR would a "CT supervision helps MR" claim be warranted --- a question we leave open rather than assert.
    \item \textbf{True external validation} on an independent multi-center cohort, beyond the internal nested CV reported here. Such a cohort with linked MR outcomes is not currently available within our consortium and is left to subsequent work.
    \item \textbf{Joint multi-task training} of the encoder with a survival auxiliary loss, the natural way to test (and potentially break) the EPV plateau.
    \item \textbf{Per-patient longitudinal modeling} across multiple scans.
\end{enumerate}


% ============================================================================
\section{Intended Use, Ethics, and Reproducibility}
\label{sec:repro}
% ============================================================================

\paragraph{Intended use.}
The pipeline is intended as clinical \emph{decision support} --- a risk-stratification aid that flags patients at elevated risk of surgery for earlier multidisciplinary review --- not an autonomous diagnostic device. The unified detector additionally enables a single deployable finding-detection model wherever either CTE or MRE is acquired.

\paragraph{Ethics.}
This retrospective multi-center study was approved by the institutional review board with a waiver of informed consent appropriate to retrospective imaging research; original imaging and clinical databases were accessed read-only. \todo{Helsinki/IRB approval numbers and the inter-center data-sharing agreements will be added at the submission stage by the clinical PI team (Shaare Zedek Medical Center).}

\paragraph{Data and code availability.}
The training, inference, and evaluation code (backbone with LoRA adapters, modality-conditioning modules, MIL head, survival fitting, decision-curve and calibration evaluation) will be released on a public GitHub repository upon acceptance of the manuscript, with the repository URL added at the proof stage. Trained checkpoints will be made available to non-commercial academic users on reasonable request. The de-identified derived dataset (10 finding probabilities and clinical features per study; no images) underlying Figures~\ref{fig:dca} and \ref{fig:km_curves} can be shared subject to the inter-center data-sharing agreements; raw DICOM data cannot be shared because original imaging is governed by individual hospital data-use agreements.


% ============================================================================
\section{Conclusion}
\label{sec:conclusion}
% ============================================================================

A single modality-conditioned foundation-model backbone can detect Crohn's-disease findings from both CT and MR enterography as well as two dedicated specialists --- halving the number of models to deploy --- \emph{provided} conditioning is lightweight. Our lead methodological result is that domain-specific normalization over-parameterizes this setting and causes statistically significant negative transfer on the dominant modality; a modality token plus FiLM is the sweet spot. Downstream, the unified model's MR-only survival head predicts time to surgery with a C-index of $0.825$ ($+0.145$ over clinical features alone), generalizes across centers, is well calibrated, and yields positive clinical net benefit. We frame Stage-1 unification as parity rather than superiority. With an updated outcomes release we extend Stage~2 to a unified CT and MR Cox model: CT delivers a meaningful prognostic signal on its own (surgery $C = 0.767$, biologic $C = 0.691$) without degrading the MR side --- an honest accounting that we believe strengthens, rather than weakens, the case for unified, conditioning-light enterography pipelines.


% ============================================================================
% FIGURES
% ============================================================================

\begin{figure}[p]
\centering
\includegraphics[width=0.98\textwidth]{figures/fig_architecture.png}
\caption{Overview of the unified two-stage pipeline and the Option~B (specialists) vs.\ Option~C (unified, modality-conditioned) comparison. CT is bowel-localized via TotalSegmentator before slice sampling and routed through the T2 branch.}
\label{fig:architecture}
\end{figure}

\begin{figure}[p]
\centering
\includegraphics[width=0.85\textwidth]{figures/fig_stage1_ladder.png}
\caption{Stage~1 macro-AUC across the conditioning ladder. Removing DSBN (A2 vs.\ A3) significantly improves MR; token+FiLM (A2) matches the MR specialist and ties the CT specialist.}
\label{fig:ladder}
\end{figure}

\begin{figure}[p]
\centering
\includegraphics[width=0.48\textwidth]{figures/calibration_5yr_surgery.png}\hfill%
\includegraphics[width=0.48\textwidth]{figures/td_auc_curve_surgery.png}
\caption{Calibration and time-dependent discrimination for surgery (unified model). Integrated Brier score $= 0.085$.}
\label{fig:calibration}
\end{figure}

\begin{figure}[p]
\centering
\includegraphics[width=0.72\textwidth]{figures/fig_dca_surgery.png}
\caption{Decision-curve analysis for surgery. The combined model has the highest net benefit across thresholds $\approx 0.03$--$0.60$.}
\label{fig:dca}
\end{figure}

\begin{figure}[p]
\centering
\includegraphics[width=0.98\textwidth]{figures/fig_km_combined.png}
\caption{Kaplan--Meier survival curves stratified by predicted-risk tertile (unified model).}
\label{fig:km_curves}
\end{figure}


% ============================================================================
% REFERENCES
% ============================================================================

\bibliographystyle{elsarticle-num}

\begin{thebibliography}{99}

\bibitem{torres2017crohns}
Torres J, Mehandru S, Colombel JF, Peyrin-Biroulet L.
Crohn's disease.
\textit{The Lancet}. 2017;389(10080):1741--1755.

\bibitem{peyrin2017lancet}
Peyrin-Biroulet L, Loftus EV, Colombel JF, Sandborn WJ.
The natural history of adult Crohn's disease in population-based cohorts.
\textit{Am J Gastroenterol}. 2010;105(2):289--297.

\bibitem{defined_mri_cd}
Rimola J, Rodriguez S, Garcia-Bosch O, et al.
Magnetic resonance for assessment of disease activity and severity in ileocolonic Crohn's disease.
\textit{Gut}. 2009;58(8):1113--1120.

\bibitem{defined_maria}
Rimola J, Ordas I, Rodriguez S, et al.
Magnetic resonance imaging for evaluation of Crohn's disease: validation of parameters of severity and quantitative index of activity.
\textit{Inflammatory Bowel Diseases}. 2011;17(8):1759--1768.

\bibitem{dl_ibd_review}
Stable E, Arnone D, Engel H, et al.
Artificial intelligence in inflammatory bowel disease: a systematic review of diagnostic and prognostic applications.
\textit{Journal of Crohn's and Colitis}. 2024;18(Supplement 1):i246--i247.

\bibitem{dinov3}
Sim\'eoni O, Vo HV, Seitzer M, Baldassarre F, Oquab M, Jose R, Khalidov V,
Szafraniec M, Yi S, Ramamonjisoa M, Massa F, Haziza D, Wehrstedt L, Wang J,
Darcet T, Moutakanni T, Sentana L, Roberts C, Vedaldi A, Tolan M, Couprie C,
Alahari K, Labatut P, Joulin A, Jegou H, Bojanowski P, Howes R.
DINOv3.
\textit{arXiv:2508.10104}. 2025.

\bibitem{baharoon2024dinov2}
Baharoon M, Qureshi W, Ouyang J, et al.
Towards general-purpose medical image representations: evaluating DINOv2 on radiology benchmarks.
\textit{arXiv:2312.02366}. 2024.

\bibitem{curia2025}
Raidium.
Curia: a multi-modal foundation model for radiology.
\textit{arXiv:2509.06830}. 2025.

\bibitem{honeybee2025}
HONeYBEE: scalable multimodal AI in oncology through foundation-model-driven embeddings.
\textit{npj Digital Medicine}. 2025. \texttt{https://doi.org/10.1038/s41746-025-02003-4}.

\bibitem{hu2021lora}
Hu EJ, Shen Y, Wallis P, et al.
LoRA: low-rank adaptation of large language models.
\textit{arXiv:2106.09685}. 2021.

\bibitem{perez2018film}
Perez E, Strub F, de Vries H, Dumoulin V, Courville A.
FiLM: visual reasoning with a general conditioning layer.
\textit{Proceedings of the AAAI Conference on Artificial Intelligence}. 2018;32(1).

\bibitem{chang2019dsbn}
Chang WG, You T, Seo S, Kwak S, Han B.
Domain-specific batch normalization for unsupervised domain adaptation.
\textit{Proceedings of the IEEE/CVF Conference on Computer Vision and Pattern Recognition (CVPR)}. 2019;7354--7362.

\bibitem{wasserthal2023totalseg}
Wasserthal J, Breit HC, Meyer MT, et al.
TotalSegmentator: robust segmentation of 104 anatomic structures in CT images.
\textit{Radiology: Artificial Intelligence}. 2023;5(5):e230024.

\bibitem{ilse2018attention}
Ilse M, Tomczak J, Welling M.
Attention-based deep multiple instance learning.
\textit{Proceedings of the 35th International Conference on Machine Learning (ICML)}. 2018;2127--2136.

\bibitem{campanella2019clinical}
Campanella G, Hanna MG, Geneslaw L, et al.
Clinical-grade computational pathology using weakly supervised deep learning on whole slide images.
\textit{Nature Medicine}. 2019;25(8):1301--1309.

\bibitem{li2020dual}
Li B, Li Y, Eliceiri KW.
Dual-stream multiple instance learning network for whole slide image classification with self-supervised contrastive learning.
\textit{Proceedings of the IEEE/CVF Conference on Computer Vision and Pattern Recognition (CVPR)}. 2021;14318--14328.

\bibitem{defined_mil_brain}
Chikontwe P, Luna M, Kang M, Hong KS, Ahn JH, Park SH.
Dual attention multiple instance learning with unsupervised complementary loss for COVID-19 screening.
\textit{Medical Image Analysis}. 2021;72:102105.

\bibitem{katzman2018deepsurv}
Katzman JL, Shaham U, Cloninger A, Bates J, Jiang T, Kluger Y.
DeepSurv: personalized treatment recommender system using a Cox proportional hazards deep neural network.
\textit{BMC Medical Research Methodology}. 2018;18(1):24.

\bibitem{lee2018deephit}
Lee C, Zame W, Yoon J, van der Schaar M.
DeepHit: a deep learning approach to survival analysis with competing risks.
\textit{Proceedings of the AAAI Conference on Artificial Intelligence}. 2018;32(1).

\bibitem{gliosurv2025}
GlioSurv: multimodal transformer for individualized survival prediction in diffuse glioma.
\textit{npj Digital Medicine}. 2025. \texttt{https://doi.org/10.1038/s41746-025-02018-x}.

\bibitem{cervipro2025}
CerviPro: deep multimodal fusion for survival prediction in locally advanced cervical cancer.
\textit{npj Digital Medicine}. 2025. \texttt{https://doi.org/10.1038/s41746-025-01903-9}.

\bibitem{prismcrc2025}
PRISM-CRC: deep multimodal fusion of patho-radiomic and clinical data for survival in colorectal cancer.
\textit{npj Digital Medicine}. 2025. \texttt{https://doi.org/10.1038/s41746-025-02210-z}.

\bibitem{ridnik2021asl}
Ridnik T, Ben-Baruch E, Zamir N, Noy A, Friedman I, Protter M, Zelnik-Manor L.
Asymmetric loss for multi-label classification.
\textit{Proceedings of the IEEE/CVF International Conference on Computer Vision (ICCV)}. 2021;82--91.

\bibitem{polsterl2020scikit}
P{\"o}lsterl S.
scikit-survival: a library for time-to-event analysis built on top of scikit-learn.
\textit{Journal of Machine Learning Research}. 2020;21(212):1--6.

\bibitem{vickers2006dca}
Vickers AJ, Elkin EB.
Decision curve analysis: a novel method for evaluating prediction models.
\textit{Medical Decision Making}. 2006;26(6):565--574.

\bibitem{defined_montreal}
Silverberg MS, Satsangi J, Ahmad T, et al.
Toward an integrated clinical, molecular and serological classification of inflammatory bowel disease.
\textit{Can J Gastroenterol}. 2005;19(Suppl A):5A--36A.

\end{thebibliography}

\end{document}
